% %--------------------------
% %	CONFIGURATION
% %--------------------------
\documentclass[11pt,a4paper]{moderncv}
\moderncvstyle{classic}
\moderncvcolor{black}
\definecolor{firstnamecolor}{rgb}{0,0,0}

% Basic packages
\usepackage[utf8]{inputenc}
\usepackage[T1]{fontenc}
\usepackage[hscale=0.7, top=2cm, bottom=2cm]{geometry}
\usepackage[dvipsnames]{xcolor}
\definecolor{DarkBlue}{RGB}{0, 0, 150}
\usepackage[colorlinks=true, urlcolor=DarkBlue, linkcolor=DarkBlue]{hyperref}

\usepackage{microtype}

% Configure font settings
\renewcommand{\rmdefault}{ppl}
\renewcommand{\sfdefault}{phv}
\renewcommand*{\namefont}{\fontsize{26}{18}\mdseries}

% Configure microtype
\microtypesetup{expansion=false}

% Load sensitive information
\input{../../secrets/personal_info_dk.tex}

% %--------------------------
\begin{document}
\makecvtitle
Hello,

\vspace*{2em}
I am writing to express my interest in contributing to Quantum Foundry Copenhagen in a technical role where my background in physics, applied machine learning, and data engineering can be useful. Your work on quantum processor manufacturing is exactly the kind of interdisciplinary engineering environment I am looking for: close to research, grounded in physical systems, and dependent on careful experimentation, automation, and reliable data.

\hspace*{2em}
My formal background is Engineering Physics from the Royal Institute of Technology (KTH), followed by a master's specialisation in Statistical Learning and Data Analytics. I graduated with a 5.0/5.0 GPA, with coursework including Quantum Physics, Solid State Physics, and Electromagnetic Theory as well as machine learning. During an exchange at Zhejiang University, I worked with a local research group on integrated photonic circuits, fabricating microring resonators and tuning their resonance frequency with controlled laser pulses to selectively remove cladding material. That work gave me hands-on experience with optical setups, lab instrumentation, device modification, and the importance of repeatable experimental protocols.

\hspace*{2em}
Professionally, I have focused on making applied ML and data systems reliable enough for real use. At Signum Life Science, I am building a Databricks forecasting platform for pharmaceutical price and demand. At Valcon, I architected an end-to-end production pipeline for a real-time neural network and led projects across pharma, energy, IoT, and manufacturing. Earlier, at RaySearch Laboratories, I worked on automated radiotherapy planning using autoencoders, sparse Gaussian processes, and optimisation on segmented 3D CT scans, in close collaboration with medical physicists and clinicians.

\hspace*{2em}
I believe this combination could be valuable in several parts of a quantum-foundry environment: Python-based automation, data acquisition and logging, analysis of experimental results, extraction of performance metrics, modelling of complex systems, and collaboration between experimental scientists, engineers, and software-oriented teams. I would also be very motivated to deepen my direct work with quantum hardware, photonics, characterization, and related laboratory workflows.

\hspace*{2em}
Outside work, I spend much of my time designing and building physical objects through woodworking and 3D printing. This reflects the same practical, curious, and creative engineering mindset I would bring to a hands-on research and engineering team.

\vspace*{2em}
Thank you for your consideration. I would be grateful for the opportunity to discuss whether my profile could be useful to Quantum Foundry Copenhagen, either now or as future roles develop. Regards,

\vspace*{1em}
Ivar Cashin Eriksson.

\end{document}
